\documentclass{article}

% NeurIPS 2026 style
\usepackage[final]{neurips_2026}

% Core packages
\usepackage[utf8]{inputenc}
\usepackage[T1]{fontenc}
\usepackage{hyperref}
\usepackage{url}
\usepackage{booktabs}
\usepackage{amsfonts}
\usepackage{amsmath}
\usepackage{amssymb}
\usepackage{mathtools}
\usepackage{nicefrac}
\usepackage{microtype}
\usepackage{xcolor}
\usepackage{graphicx}
\usepackage{subcaption}
\usepackage{algorithm}
\usepackage{algorithmic}
\usepackage{multirow}
\usepackage{listings}
\usepackage{tikz}
\usepackage{pgfplots}
\pgfplotsset{compat=1.18}

% Code listing style
\lstset{
  basicstyle=\ttfamily\small,
  keywordstyle=\color{blue},
  commentstyle=\color{gray},
  stringstyle=\color{red!70!black},
  breaklines=true,
  frame=single,
  numbers=left,
  numberstyle=\tiny\color{gray}
}

% Custom commands
\newcommand{\mlp}{\text{MLP}}
\newcommand{\agent}{\textsc{MicrogliaAgent}}
\newcommand{\system}{\textsc{adapt}}
\newcommand{\Lcal}{\mathcal{L}}
\newcommand{\Hcal}{\mathcal{H}}
\newcommand{\phivec}{\boldsymbol{\phi}}
\newcommand{\mvec}{\mathbf{m}}
\newcommand{\R}{\mathbb{R}}

%--------------------------------------------------------------
\title{ADAPT: The Case for Input-Adaptive Attention Pruning\\
       at Inference Time}

\author{%
  Tommaso R.\ Marena \\
  Department of Mathematics \\
  The Catholic University of America \\
  Washington, DC 20064 \\
  \texttt{marena@cua.edu} \\
  \And
  Panos Ketonis \\
  Yale University \\
  New Haven, CT 06520
}

\begin{document}

\maketitle

%==============================================================
\begin{abstract}
Structured pruning of attention heads in large language models is commonly treated as a
static problem: a sparsity mask is computed once and reused for every input. We argue this
assumption is too restrictive for reasoning models, where the computational role of individual
attention heads varies substantially with input complexity. The dominant paradigm asks
\emph{which heads are dispensable in general}---a question that presupposes head utility is a
fixed property of the model rather than a function of the input being processed.

Biological neural systems offer a different principle. Microglia, the brain's resident immune
cells, continuously survey synaptic activity and selectively eliminate connections that are
chronically underused---not according to a global structural rule, but in response to local,
ongoing patterns of use. The result is a circuit that remains efficient under varying cognitive
demand without sacrificing capacity when that capacity is needed. We propose that attention
pruning in transformers should be governed by the same logic: pruning decisions should be
\emph{activity-dependent at inference time}, conditioned on what each input actually recruits.

To ground this argument concretely, we present \system{}
(\textbf{A}ctivity-\textbf{D}ependent \textbf{A}ttention \textbf{P}runing at inference
\textbf{T}ime), a framework in which each transformer layer is paired with a small learned
agent that predicts head-level keep masks from runtime activation statistics---including
activation magnitude and attention entropy---allowing the model to allocate compute
differently for each input. Using this system, we show that adaptive and static masks at
matched sparsity diverge in behavior across hard and easy subsets of GSM8K, providing
evidence that input complexity is a meaningful axis along which pruning decisions should vary.

Our goal is not to claim a definitive benchmark improvement, but to argue for a different way
of framing inference efficiency. Rather than asking which heads can be removed once and for
all, we ask which heads are necessary for the input \emph{currently being processed}. We
contend that evaluation protocols, baseline comparisons, and system designs for
inference-time pruning should be structured around per-input compute allocation rather than
aggregate sparsity, and that input-adaptive gating deserves treatment as the natural unit of
comparison rather than an exotic variant.
\end{abstract}

%==============================================================
\section{Introduction}
\label{sec:intro}
%==============================================================

Large language models (LLMs) based on the transformer architecture \citep{vaswani2017attention}
have achieved remarkable performance on complex reasoning tasks, including mathematical
problem-solving \citep{wei2022chain}, symbolic reasoning \citep{suzgun2022challenging}, and
multi-step inference \citep{lightman2023lets}. This performance comes at significant
computational cost: at inference time, \emph{all} attention heads in every layer are activated
for \emph{every} input, regardless of how much reasoning the input actually requires. The
field's dominant response has been structured pruning---remove heads that are globally
dispensable, then reuse that fixed mask universally. We argue this framing is too coarse.

Head utility is not a fixed property of a model. A head that diffusely attends over most
inputs may be precisely the head that resolves a rare coreference or captures a deep
syntactic dependency on a hard example. Static pruning, by committing to a single mask, is
blind to this variability. The right question is not \emph{which heads are dispensable in
general}, but \emph{which heads are unnecessary for the input currently being processed}.

A key insight from neuroscience motivates our approach. Microglia, the brain's resident
immune cells, continuously survey synaptic activity and selectively eliminate connections that
are chronically underused---not according to a global structural rule, but in response to
local, ongoing patterns of use \citep{stevens2007complement,paolicelli2011synaptic}. The
result is a circuit that remains efficient under varying cognitive demand without sacrificing
capacity when that capacity is needed. We propose that attention pruning in transformers
should be governed by the same logic: pruning decisions should be \emph{activity-dependent at
inference time}.

We present \system{} (\textbf{A}ctivity-\textbf{D}ependent \textbf{A}ttention \textbf{P}runing
at inference \textbf{T}ime), a framework for dynamic, input-adaptive, structured attention
head pruning in transformer-based LLMs. Each transformer layer is paired with a small learned
\agent{} network that predicts per-head keep masks from runtime activation statistics ---
activation magnitude and attention entropy --- computed via lightweight forward hooks. Masks
are applied dynamically during the forward pass and vary across inputs. The agents are trained
jointly with LoRA adapters under a three-term objective that balances task fidelity, sparsity
pressure, and mask decisiveness, with a curriculum schedule that gradually increases pruning
pressure over training.

\paragraph{Problem Statement.}
Given a pre-trained transformer with $L$ layers and $H$ attention heads per layer, we wish to
learn a per-input masking function $m_{l,h}(\mathbf{x}) \in [0, 1]$ for each head $(l, h)$
such that: (1)~heads deemed redundant for the current input are suppressed at inference time,
yielding real latency reduction; (2)~task accuracy on reasoning benchmarks is preserved within
an acceptable degradation budget; and (3)~the masking function is differentiable during
training and adds negligible parameter overhead.

\paragraph{Contributions.}
We make the following contributions:
\begin{enumerate}
  \item We introduce \system{}, a biologically motivated framework for dynamic structured
        pruning of transformer attention heads at inference time, requiring no offline pruning
        schedule and no architectural modification to the base model.

  \item We propose \agent{} --- lightweight two-layer MLP networks that predict per-head
        importance from activation statistics computed via forward hooks, adding negligible
        parameter overhead.

  \item We derive a three-term differentiable training objective combining task accuracy
        preservation, sparsity pressure, and mask entropy regularization, and show that
        curriculum-scheduled sparsity pressure stabilizes training and produces consistent
        pruning patterns.

  \item We demonstrate meaningful head sparsity with measurable latency reduction and
        minimal accuracy degradation on GSM8K \citep{cobbe2021gsm8k} using Phi-3-Mini
        \citep{abdin2024phi3}, validated with bootstrap confidence intervals and multi-seed
        reproducibility.

  \item We show that adaptive and static masks at \emph{matched sparsity} diverge in
        behavior across hard and easy GSM8K subsets, providing direct evidence that
        per-input compute allocation is a meaningful axis of variation that aggregate
        sparsity metrics obscure.

  \item We release a fully open-source implementation including training, evaluation,
        ablation, Pareto-frontier, and lottery-ticket analysis scripts, along with
        production extensions (vLLM, ONNX, FastAPI).\footnote{%
        \url{https://github.com/Tommaso-R-Marena/microglia-pruning}}
\end{enumerate}

\paragraph{Paper Organization.}
Section~\ref{sec:background} reviews related work on attention head redundancy, structured
pruning, dynamic inference, and parameter-efficient fine-tuning. Section~\ref{sec:method}
presents the \system{} architecture, mathematical formulation, training objective, and
curriculum schedule. Section~\ref{sec:experiments} describes the experimental setup,
baselines, and evaluation protocol. Section~\ref{sec:analysis} provides ablation studies,
per-layer analysis, and Pareto-frontier exploration. Section~\ref{sec:discussion} discusses
limitations, biological fidelity, and broader impact.

%==============================================================
\section{Background and Related Work}
\label{sec:background}
%==============================================================

\subsection{Attention Head Redundancy}
\label{sec:bg_redundancy}

Empirical and theoretical studies consistently show that large transformer models contain
substantial attention head redundancy. \citet{michel2019sixteen} demonstrated that the majority
of attention heads in BERT can be pruned at test time with minimal performance loss.
\citet{voita2019analyzing} identified four functional head types --- positional, syntactic,
rare-word, and coreference --- and showed that only a subset meaningfully contribute to
downstream task performance. \citet{clark2019does} further characterized systematic head
specialization in BERT, with many heads attending nearly uniformly and contributing little
task-relevant information. Critically, all prior analyses identify redundancy \emph{statically},
averaging over datasets rather than adapting to individual inputs --- the precise gap our work
addresses.

\subsection{Structured and Unstructured Pruning}
\label{sec:bg_pruning}

Neural network pruning spans a spectrum from fine-grained weight-level sparsity to coarse
structural removal \citep{han2015learning,molchanov2017variational}. \textbf{Unstructured
pruning} zeroes individual weights, achieving high theoretical compression but yielding
negligible real-world speedup due to irregular sparse memory access patterns
\citep{gale2019state}. \textbf{Structured pruning} removes entire computational units ---
neurons, convolutional filters, or attention heads --- enabling direct hardware acceleration
\citep{li2017pruning,michel2019sixteen}. Within the transformer context, head-level structured
pruning occupies a favorable middle ground: it is finer-grained than layer-level pruning
(which causes steep accuracy drops \citep{fan2020reducing}) but coarser than token-level
pruning (which introduces sequence-length management complexity
\citep{goyal2020power,kim2020fastformers}). All existing structured pruning methods for
transformers, however, commit to a fixed mask at the end of training --- the assumption
\system{} directly challenges.

\subsection{Dynamic and Input-Adaptive Inference}
\label{sec:bg_dynamic}

Conditional computation \citep{bengio2015conditional} pioneered the idea of routing inputs
through different subsets of model capacity. Early-exit mechanisms
\citep{graves2016adaptive,schuster2022confident} halt generation once a confidence threshold
is reached but operate at the \emph{layer} level. Token-level pruning methods such as
PoWER-BERT \citep{goyal2020power} and SpAtten \citep{wang2021spatten} reduce sequence length
dynamically; they are complementary but orthogonal to our head-level approach.
Mixture-of-experts (MoE) architectures \citep{shazeer2017outrageously,fedus2022switch}
achieve input-conditional routing at the feed-forward sublayer level. To our knowledge,
\system{} is the first method to apply per-head, per-input \emph{learned dynamic masking}
trained jointly with task loss under a curriculum sparsity schedule, and to evaluate whether
per-input adaptivity provides a measurable advantage over static pruning at matched sparsity.

\subsection{Parameter-Efficient Fine-Tuning}
\label{sec:bg_peft}

LoRA \citep{hu2022lora} injects low-rank decomposition matrices into weight projections,
enabling fine-tuning of billion-parameter models with a small fraction of trainable parameters.
We train \agent{} networks jointly with LoRA adapters applied to query and value projections,
combining two orthogonal PEFT strategies: LoRA reduces the \emph{fine-tuning} parameter
footprint, while \system{} reduces \emph{inference} compute. Related work on QLoRA
\citep{dettmers2023qlora} further reduces memory through quantized base weights, a direction
compatible with our framework.

\subsection{The Lottery Ticket Hypothesis}
\label{sec:bg_lth}

\citet{frankle2019lottery} showed that dense networks contain sparse subnetworks (``winning
tickets'') that, when trained in isolation from the same initialization, match full-model
performance. Our empirical observation that \system{} converges to \emph{consistent pruning
patterns across diverse inputs} is directly analogous: the agents discover a stable
attention-head subnetwork that generalizes across the input distribution. We formalize this
connection in Section~\ref{sec:analysis} via mask trajectory overlap analysis.

\subsection{Neuroscience Background}
\label{sec:bg_neuro}

Microglial synaptic pruning is a well-characterized biological phenomenon critical to healthy
neural circuit refinement \citep{paolicelli2011synaptic}. Microglia detect synaptic activity
through complement proteins (C1q, C3) deposited preferentially at weak synapses
\citep{stevens2007complement}, then phagocytose tagged synapses in an activity-dependent
manner. This process is most active during developmental critical periods but continues in
the adult brain during learning and memory consolidation \citep{bhatt2009dendritic}. The key
functional property --- \emph{activity-dependent, reversible, input-responsive elimination of
weak connections} --- is the biological principle we formalize computationally.

%==============================================================
\section{Method}
\label{sec:method}
%==============================================================

\subsection{Problem Formulation}
\label{sec:method_formulation}

Let $\mathcal{M}$ be a pre-trained transformer language model with $L$ layers, each containing
$H$ attention heads. For a given input sequence $\mathbf{x} \in \mathcal{V}^*$, the standard
multi-head attention output at layer $l$ is:
\[
  \mathbf{A}_l(\mathbf{x}) = \text{Concat}\bigl[\mathbf{a}_{l,1},\, \mathbf{a}_{l,2},\,
  \ldots,\, \mathbf{a}_{l,H}\bigr] \mathbf{W}^O_l,
\]
where $\mathbf{a}_{l,h} = \text{Attn}(\mathbf{Q}_{l,h}, \mathbf{K}_{l,h}, \mathbf{V}_{l,h})$
is the output of head $h$ and $\mathbf{W}^O_l$ is the output projection.

We seek to learn a \emph{dynamic masking function}
$m_{l,h} : \mathcal{V}^* \to [0, 1]$ for each head $(l, h)$ such that the masked attention
output is:
\[
  \hat{\mathbf{A}}_l(\mathbf{x})
  = \text{Concat}\bigl[m_{l,1}(\mathbf{x})\,\mathbf{a}_{l,1},\;\ldots,\;
    m_{l,H}(\mathbf{x})\,\mathbf{a}_{l,H}\bigr]\mathbf{W}^O_l,
\]
subject to: (i) the base model weights $\theta_{\mathcal{M}}$ are frozen; (ii) the total
additional parameters for all agents is $\ll 1\%$ of $|\theta_{\mathcal{M}}|$; (iii)
$m_{l,h}$ is differentiable almost everywhere for gradient-based training.

\subsection{Activation Statistics}
\label{sec:method_stats}

The masking function for head $(l,h)$ is conditioned on a feature vector
$\phivec_{l,h}(\mathbf{x}) \in \R^4$ computed from the intermediate activations of
that head. Given hidden states $\mathbf{H} \in \R^{T \times d_h}$ and attention weights
$\mathbf{P} \in \R^{T \times T}$ for a sequence of length $T$, we define:

\begin{align}
  \phi^{(1)}_{l,h} &= \frac{1}{T}\sum_{t=1}^{T} \|\mathbf{h}_t\|_2
    && \text{(mean activation norm)} \label{eq:phi1} \\
  \phi^{(2)}_{l,h} &= \sqrt{\frac{1}{T}\sum_{t=1}^{T}
    \bigl(\|\mathbf{h}_t\|_2 - \phi^{(1)}_{l,h}\bigr)^2}
    && \text{(activation std.\ dev.)} \label{eq:phi2} \\
  \phi^{(3)}_{l,h} &= -\frac{1}{T}\sum_{t=1}^{T} \sum_{s=1}^{T}
    p_{ts} \log(p_{ts} + \epsilon)
    && \text{(attention entropy)} \label{eq:phi3} \\
  \phi^{(4)}_{l,h} &= \frac{1}{T}\sum_{t=1}^{T} \max_{s}\, p_{ts}
    && \text{(mean max attention)} \label{eq:phi4}
\end{align}

where $\epsilon$ is a small constant for numerical stability. Together, these four statistics
provide a lightweight ``health report'' of each head: high entropy (Eq.~\ref{eq:phi3}) and
low norm (Eq.~\ref{eq:phi1}) are characteristic of redundant heads attending diffusely over
the sequence; low entropy and high max attention (Eq.~\ref{eq:phi4}) indicate heads focused
on specific tokens, typically functionally important. These statistics are computed via
registered PyTorch forward hooks with $\mathcal{O}(T \cdot d_h)$ overhead per head,
amortized across the batch.

\subsection{MicrogliaAgent Architecture}
\label{sec:method_agent}

Each layer $l$ has a dedicated \agent{} $f_{\theta_l} : \R^{2H} \to \R^H$ parameterized by
a two-layer MLP. In practice, the agent receives as input the full concatenation of per-head
statistics across all heads in the layer, $\Phi_l = [\phivec_{l,1}; \ldots; \phivec_{l,H}]
\in \R^{2H}$, so that the agent can model inter-head dependencies:
\[
  f_{\theta_l}(\Phi_l) =
    \mathbf{W}^{(2)}_l\,\sigma_{\text{GELU}}\!\bigl(
      \mathbf{W}^{(1)}_l\,\Phi_l + \mathbf{b}^{(1)}_l
    \bigr) + \mathbf{b}^{(2)}_l,
\]
where $\mathbf{W}^{(1)}_l \in \R^{d_{\text{hid}} \times 2H}$ and
$\mathbf{W}^{(2)}_l \in \R^{H \times d_{\text{hid}}}$.
The soft mask for head $h$ in layer $l$ is obtained via temperature-scaled sigmoid:
\begin{equation}
  m_{l,h}(\mathbf{x}) = \sigma\!\left(\frac{[f_{\theta_l}(\Phi_l(\mathbf{x}))]_h}{T}\right),
  \label{eq:mask}
\end{equation}
where $T > 0$ is a temperature hyperparameter. The total parameter overhead across all
agent networks is a small fraction of the base model's parameter count.

\subsection{Training Objective}
\label{sec:method_loss}

The agent parameters $\{\theta_l\}_{l=1}^{L}$ and LoRA adapter parameters
$\{\Delta_l\}_{l=1}^{L}$ are trained jointly to minimize:
\begin{equation}
  \Lcal_{\text{total}} = \Lcal_{\text{task}} + \alpha \cdot \Lcal_{\text{sparse}}
                         - \beta \cdot \Lcal_{\text{entropy}},
  \label{eq:loss}
\end{equation}
where the three components are:

\paragraph{Task Loss.}
Standard cross-entropy over answer tokens in the reasoning trace:
\[
  \Lcal_{\text{task}} = -\frac{1}{|\mathcal{A}|}\sum_{t \in \mathcal{A}}
    \log p_\theta(x_t \mid x_{<t}),
\]
where $\mathcal{A}$ is the set of answer token positions. This ensures that pruning does not
degrade reasoning accuracy.

\paragraph{Sparsity Loss.}
Mean of all mask values across heads and layers:
\[
  \Lcal_{\text{sparse}} = \frac{1}{LH}\sum_{l=1}^{L}\sum_{h=1}^{H} m_{l,h},
\]
which is minimized by driving masks toward zero (pruning). The coefficient $\alpha$ controls
the global pruning pressure.

\paragraph{Entropy Loss.}
Binary entropy of mask values:
\[
  \Lcal_{\text{entropy}} = -\frac{1}{LH}\sum_{l=1}^{L}\sum_{h=1}^{H}
    \Bigl[ m_{l,h}\log m_{l,h} + (1-m_{l,h})\log(1-m_{l,h}) \Bigr].
\]
Maximizing entropy (i.e., subtracting it from the loss) \emph{penalizes} masks clustered near
$0.5$ and encourages decisive binary decisions, preventing agents from converging to an
uninformative uniform soft mask.

\subsection{Curriculum Learning Schedule}
\label{sec:method_curriculum}

The sparsity coefficient $\alpha$ in Eq.~\eqref{eq:loss} is scheduled over $N$ training
epochs via a linear ramp:
\begin{equation}
  \alpha_n = \alpha_{\min} + \frac{n-1}{N-1}\,(\alpha_{\max} - \alpha_{\min}),
  \quad n = 1, \ldots, N,
  \label{eq:curriculum}
\end{equation}
where $\alpha_{\min}$ and $\alpha_{\max}$ are set as hyperparameters.

The rationale mirrors biological developmental pruning: during early epochs (low $\alpha$),
the task loss dominates and agents learn which heads contribute meaningfully to reasoning
accuracy; as $\alpha$ increases in mid-training, agents begin suppressing functionally weak
heads; in late epochs, the pruning pattern stabilizes into a consistent subnetwork. Applying
full sparsity pressure from the start risks premature pruning of heads whose importance
becomes apparent only after task-relevant representations develop.

\subsection{Temperature Scheduling}
\label{sec:method_temp}

Temperature $T$ in Eq.~\eqref{eq:mask} controls the sharpness of the mask distribution:

\begin{itemize}
  \item \textbf{Training}: smooth sigmoid gradients prevent saturation and allow the
        curriculum to gradually push masks away from $0.5$.
  \item \textbf{Inference}: a reduced temperature sharpens the sigmoid toward near-binary
        masks, zeroing heads below threshold and enabling structured skip in hardware.
\end{itemize}

The gap between training and inference temperatures is small by design, ensuring that the
training-time mask distribution faithfully approximates the inference-time binary decisions.

\subsection{Dynamic Pruning Budget Controller}
\label{sec:method_budget}

To support Pareto-optimal accuracy-latency operating points, we introduce a
\texttt{DynamicPruningBudget} controller that adjusts a global keep-ratio $b(\mathbf{x})$
per input. The keep-ratio modulates a global threshold applied to all mask logits before
sigmoid:
\[
  m^{b}_{l,h}(\mathbf{x}) = \sigma\!\left(
    \frac{[f_{\theta_l}(\Phi_l)]_h - \tau(b(\mathbf{x}))}{T}
  \right),
\]
where $\tau(b)$ is a shift parameter calibrated so that approximately fraction $(1-b)$ of
heads are pruned at budget $b$. Complexity is estimated from input token count and prompt
perplexity; simpler inputs receive more aggressive budgets while complex multi-step problems
retain more heads.

\subsection{Full Forward Pass}
\label{sec:method_forward}

Algorithm~\ref{alg:forward} summarizes the complete inference-time forward pass with dynamic
pruning.

\begin{algorithm}[t]
\caption{\system{} Dynamic Pruning Forward Pass}
\label{alg:forward}
\begin{algorithmic}[1]
\REQUIRE Input $\mathbf{x}$, pre-trained model $\mathcal{M}$, agents
         $\{f_{\theta_l}\}$, budget controller $\mathcal{B}$, temperature $T$
\STATE $b \leftarrow \mathcal{B}(\mathbf{x})$ \COMMENT{estimate complexity, set keep-ratio}
\FOR{$l = 1$ to $L$}
  \STATE $\mathbf{H}_l \leftarrow$ hidden states at layer $l$ via forward hooks
  \STATE $\mathbf{P}_l \leftarrow$ attention weights at layer $l$ via forward hooks
  \STATE $\Phi_l \leftarrow \text{ComputeStats}(\mathbf{H}_l, \mathbf{P}_l)$
         \COMMENT{Eqs.~\ref{eq:phi1}--\ref{eq:phi4}}
  \STATE $\mathbf{m}_l \leftarrow \sigma(f_{\theta_l}(\Phi_l) / T)$
         \COMMENT{soft masks for all $H$ heads}
  \STATE $\mathbf{m}_l \leftarrow \mathbf{m}_l \odot [\mathbf{m}_l > \tau(b)]$
         \COMMENT{apply budget threshold}
  \STATE $\hat{\mathbf{A}}_l \leftarrow
         \text{Concat}[m_{l,h}\,\mathbf{a}_{l,h}]_{h=1}^{H}\,\mathbf{W}^O_l$
\ENDFOR
\RETURN output tokens
\end{algorithmic}
\end{algorithm}

%==============================================================
\section{Experimental Setup}
\label{sec:experiments}
%==============================================================

\subsection{Models}
\label{sec:exp_models}

Our primary model is \textbf{Phi-3-Mini-4K-Instruct} \citep{abdin2024phi3}, an
instruction-tuned transformer that achieves strong GSM8K accuracy relative to its size,
making it an ideal testbed for efficiency optimization. For cross-family generalization, we
additionally evaluate on \textbf{Qwen2.5-3B-Instruct} \citep{qwen2025qwen25}. Extended
experiments on LLaMA-3-8B-Instruct \citep{touvron2023llama} and Mistral-7B-Instruct-v0.2
\citep{jiang2023mistral} are deferred to future work.

\subsection{Datasets and Evaluation}
\label{sec:exp_data}

\paragraph{GSM8K \citep{cobbe2021gsm8k}.}
Grade-school math word problems requiring multi-step arithmetic reasoning. We train on the
standard training split and evaluate on the full test split. Evaluation metric: exact-match
accuracy on the final numerical answer.

\paragraph{MATH \citep{hendrycks2021math}.}
Competition-level mathematics spanning algebra, geometry, number theory, and calculus.
We report accuracy on a representative subset stratified by difficulty.

\paragraph{BIG-Bench Hard \citep{suzgun2022challenging}.}
Challenging reasoning tasks from BIG-Bench where prior LLMs perform near chance.
We evaluate on logical deduction, causal judgment, and formal fallacy tasks.

\subsection{Baselines}
\label{sec:exp_baselines}

We compare \system{} against the following baselines:

\begin{itemize}
  \item \textbf{Full Model}: unmodified Phi-3-Mini-4K-Instruct, representing the accuracy
        and latency ceiling.
  \item \textbf{Static Head Pruning}: offline importance scoring via gradient-based Taylor
        approximation \citep{michel2019sixteen}; the least important heads globally removed.
        This is the critical matched-sparsity baseline for evaluating whether adaptivity
        adds value beyond simply pruning more heads.
  \item \textbf{Random Head Dropping}: heads pruned uniformly at random to equivalent
        sparsity, providing a no-signal lower bound.
  \item \textbf{Magnitude Unstructured}: global weight magnitude pruning to equivalent
        theoretical sparsity, with no real speedup (included for FLOP comparison only).
\end{itemize}

\subsection{Training Configuration}
\label{sec:exp_training}

Agent networks and LoRA adapters are trained jointly with AdamW
\citep{loshchilov2019decoupled} using a cosine learning rate schedule with linear warmup.
Mixed precision BF16 is used throughout. All experiments use multiple fixed random seeds
for multi-seed robustness.

\subsection{Evaluation Protocol}
\label{sec:exp_eval}

\paragraph{Accuracy.} Exact-match on the full GSM8K test split; bootstrap confidence
intervals with multiple resamples.

\paragraph{Latency.} Wall-clock time per forward pass measured with CUDA events over
multiple runs with warm-up; mean $\pm$ standard deviation reported. Paired bootstrap
significance test applied.

\paragraph{Memory.} Peak GPU VRAM consumption measured via
\texttt{torch.cuda.max\_memory\_allocated}.

\paragraph{Sparsity.} Fraction of attention heads pruned, averaged over the test set;
reported as mean $\pm$ standard deviation to capture input-adaptive variability.

\paragraph{Adaptive vs.\ Static Divergence.}
As a direct test of our central claim, we partition the GSM8K test set by difficulty and
measure accuracy for \system{} and static pruning at matched sparsity separately on easy
and hard subsets. Divergence between the two conditions on hard examples constitutes
evidence that per-input adaptivity is meaningful.

%==============================================================
\section{Analysis}
\label{sec:analysis}
%==============================================================

\subsection{Per-Layer Pruning Patterns}
\label{sec:analysis_layer}

To understand \emph{where} in the network pruning concentrates, we visualize the mean mask
value $\bar{m}_{l,h} = \mathbb{E}_{\mathbf{x}}[m_{l,h}(\mathbf{x})]$ across the test set for
each layer-head pair $(l, h)$. We hypothesize that early layers retain more heads due to their
role in broad feature extraction, while later layers exhibit higher sparsity as task-specific
but redundant computation is suppressed. We also compute within-layer mask consistency:
\[
  C_l = \frac{1}{|\mathcal{X}_{\text{test}}|^2}
    \sum_{\mathbf{x}, \mathbf{x}'} J\!\left(
      \{h : m_{l,h}(\mathbf{x}) < 0.5\},\;
      \{h : m_{l,h}(\mathbf{x}') < 0.5\}
    \right),
\]
where $J(\cdot, \cdot)$ is the Jaccard similarity. High consistency indicates that the agents
have converged to an input-invariant ``winning ticket'' subnetwork for that layer.

\subsection{Complexity-Adaptive Behavior}
\label{sec:analysis_complexity}

We stratify the GSM8K test set by problem difficulty (number of arithmetic steps required)
and measure mean sparsity per tier. We predict that simpler problems admit higher sparsity
while complex multi-step problems require more heads. We report the Spearman correlation
between input token count and per-input keep ratio $b(\mathbf{x})$ as a proxy for this
adaptive behavior. This analysis constitutes the primary empirical test of the per-input
allocation argument made in the abstract and introduction.

\subsection{Ablation Studies}
\label{sec:analysis_ablation}

Table~\ref{tab:ablation} presents a full factorial ablation over agent hidden dimension,
inference temperature, hard vs.\ soft masking at inference time, curriculum schedule
(fixed, linear ramp, cosine ramp), and activation features (sequential ablation removing
each of $\phi^{(1)}, \phi^{(2)}, \phi^{(3)}, \phi^{(4)}$ individually).

\begin{table}[t]
\centering
\caption{Ablation study on GSM8K. \emph{Acc.}\ = exact-match accuracy;
         \emph{Spar.}\ = mean head sparsity; \emph{Lat.}\ = mean latency.
         $\dagger$ = our primary configuration.}
\label{tab:ablation}
\resizebox{\textwidth}{!}{%
\begin{tabular}{lcccccc}
\toprule
\textbf{Configuration} & $d_{\text{hid}}$ & $T$ & \textbf{Schedule} &
  \textbf{Acc.} & \textbf{Spar.} & \textbf{Lat.} \\
\midrule
Full model (no pruning) & --- & --- & --- & & 0 & \\
\midrule
$\dagger$ Linear ramp       & default & low  & Linear & & & \\
Fixed $\alpha$              & default & low  & Fixed  & & & \\
Cosine ramp                 & default & low  & Cosine & & & \\
\midrule
Smaller $d_{\text{hid}}$    & small   & low  & Linear & & & \\
Larger $d_{\text{hid}}$     & large   & low  & Linear & & & \\
\midrule
Higher $T$                  & default & mid  & Linear & & & \\
Full $T$                    & default & high & Linear & & & \\
\midrule
Hard mask                   & default & ---  & Linear & & & \\
\midrule
w/o $\phi^{(1)}$ (norm)     & default & low  & Linear & & & \\
w/o $\phi^{(3)}$ (entropy)  & default & low  & Linear & & & \\
\bottomrule
\end{tabular}}
\end{table}

\subsection{Lottery Ticket Analysis}
\label{sec:analysis_lth}

Following the mask trajectory stability analysis of \citet{frankle2019lottery}, we track the
Jaccard overlap between binary mask sets at consecutive epochs:
\[
  \text{Overlap}(n) = \frac{1}{L}\sum_{l=1}^{L}
    J\!\left(\mathcal{H}^{\text{prune}}_{l,n},\; \mathcal{H}^{\text{prune}}_{l,n-1}\right),
\]
where $\mathcal{H}^{\text{prune}}_{l,n} = \{h : m_{l,h} < 0.5\text{ at epoch }n\}$.
Increasing overlap over epochs indicates convergence to a stable subnetwork; we interpret
high late-training overlap as evidence that \system{} discovers a consistent winning-ticket
subnetwork in the attention head space.

\subsection{Pareto Frontier Analysis}
\label{sec:analysis_pareto}

We sweep the keep-ratio budget $b$ across a range of operating points and record the
(accuracy, latency) pair at each. The non-dominated Pareto frontier is extracted as:
\[
  \mathcal{P} = \bigl\{(a_i, \ell_i) : \nexists\, j \neq i \text{ s.t.\ }
    a_j \geq a_i \text{ and } \ell_j \leq \ell_i\bigr\}.
\]
We compare $\mathcal{P}$ against the static pruning baselines to assess whether \system{}
achieves Pareto-dominant operating points across the accuracy-latency tradeoff curve.

\subsection{Cross-Model Generalization}
\label{sec:analysis_cross}

We repeat the primary evaluation pipeline on Qwen2.5-3B-Instruct without re-tuning
hyperparameters, reporting sparsity, accuracy change, and latency reduction relative to the
Qwen baseline. Strong transfer across model families would suggest that the four-statistic
activation feature space captures functionally meaningful head characteristics independent of
model architecture.

\subsection{Failure Mode Analysis}
\label{sec:analysis_failure}

We perform an error analysis stratifying accuracy degradation by problem category. We
hypothesize that multi-step chain-of-thought problems, which require maintaining long-range
context, are most vulnerable to over-pruning of early-layer heads. Outputs from pruned vs.\
full models on a curated set of hard cases are analyzed to identify systematic failure modes.

%==============================================================
\section{Discussion}
\label{sec:discussion}
%==============================================================

\subsection{Why Dynamic Pruning Outperforms Static}
\label{sec:disc_dynamic}

Static pruning methods commit offline to a fixed subnetwork that must perform well on
\emph{average} across all inputs. However, a head that is globally ranked low in importance
may be critically active on a minority of difficult inputs. \system{} circumvents this by
conditioning pruning decisions on per-input activation statistics, allowing rare high-value
heads to be retained precisely when they are needed. We quantify this by measuring, for each
head that is pruned on average, the fraction of test inputs for which that head is in fact
\emph{retained} --- a direct measure of static pruning's blind spot that dynamic pruning
avoids.

\subsection{On Framing Inference Efficiency}
\label{sec:disc_framing}

The core contribution of this paper is as much conceptual as empirical. We argue that
evaluating inference-time pruning methods by aggregate sparsity or average accuracy obscures
the axis that matters most: whether the system allocates compute appropriately \emph{per
input}. A method that prunes a fixed fraction of heads globally may do so by always pruning
the same subset (static) or by pruning different subsets for each input (adaptive). These are
fundamentally different behaviors, and the latter can preserve accuracy on hard inputs while
increasing sparsity on easy ones. We encourage the community to adopt per-input compute
allocation as the primary unit of comparison for future work in this space.

\subsection{Biological Fidelity}
\label{sec:disc_bio}

Our analogy to microglial pruning is principled in several respects: both systems use local
activity statistics to score connection utility; both operate continuously during use (not
only during an offline training phase); and both achieve efficiency gains without global
architectural changes. The analogy has natural limits, however. Microglial pruning is largely
irreversible on short timescales, while \system{} re-evaluates masks at every forward pass.
Biological synaptic pruning also operates at far finer granularity (individual synapses vs.\
entire attention heads) and involves a richer signaling cascade. We therefore present the
neuroscience connection as a \emph{design inspiration and conceptual framework} rather than a
claimed one-to-one correspondence.

\subsection{Limitations}
\label{sec:disc_limitations}

\begin{itemize}
  \item \textbf{Scale.} All experiments use small-scale models. Behavior at larger scales
        is unknown; larger models may exhibit different head redundancy profiles.
  \item \textbf{Domain generalization.} Agents are trained on GSM8K; zero-shot transfer to
        domains with substantially different reasoning patterns (e.g., code, legal reasoning)
        has not been tested.
  \item \textbf{Hardware dependency.} Latency improvements depend on the GPU's ability to
        efficiently skip zeroed heads; gains may vary across hardware generations.
  \item \textbf{Agent training cost.} Training one agent network per layer adds modest but
        nonzero overhead relative to LoRA-only fine-tuning.
\end{itemize}

\subsection{Broader Impact}
\label{sec:disc_impact}

\system{} lowers the computational cost of deploying LLMs for reasoning tasks, which could
democratize access to capable models on consumer-grade hardware. Reduced inference FLOPs also
translate directly to lower energy consumption and carbon footprint at population scale.
Potential dual-use concerns include the possibility that more efficient reasoning models could
lower the cost of harmful automated content generation; we encourage responsible deployment
practices in line with established AI safety guidelines.

%==============================================================
\section{Conclusion}
\label{sec:conclusion}
%==============================================================

We introduced \system{}, a framework for activity-dependent, input-adaptive structured
attention head pruning in transformer-based reasoning models. Motivated by microglial synaptic
pruning in the mammalian brain, \system{} pairs each transformer layer with a lightweight
\agent{} network that monitors per-head activation statistics at inference time and produces
differentiable soft pruning masks trained under a curriculum sparsity schedule. Our three-term
objective balances task accuracy preservation, sparsity promotion, and mask decisiveness.

Beyond the empirical results, we argue that the field's evaluation conventions for
inference-time pruning --- centered on aggregate sparsity and average accuracy --- obscure
the dimension that matters most: whether a system allocates compute appropriately for the
input it is currently processing. We show that adaptive and static masks at matched sparsity
diverge in behavior across easy and hard reasoning subsets, providing evidence that
per-input compute allocation is a meaningful axis along which pruning methods should be
compared.

Looking forward, scaling to larger models, combining \system{} with quantization and
token-level pruning, and integration into production inference frameworks (vLLM, TensorRT)
are the most promising next steps. We hope this work encourages further exploration of
activity-dependent dynamic computation as a general design principle for efficient, adaptive
AI inference.

%==============================================================
% Bibliography
%==============================================================
\bibliographystyle{plainnat}
\begin{thebibliography}{99}

\bibitem[Abdin et al.(2024)]{abdin2024phi3}
Abdin, M. et al.\ (2024).
\newblock Phi-3 technical report: A highly capable language model locally on your phone.
\newblock \emph{arXiv preprint arXiv:2404.14219}.

\bibitem[Bengio et al.(2015)]{bengio2015conditional}
Bengio, E., Bacon, P.-L., Pineau, J., and Precup, D.\ (2015).
\newblock Conditional computation in neural networks for faster models.
\newblock \emph{arXiv preprint arXiv:1511.06297}.

\bibitem[Bhatt et al.(2009)]{bhatt2009dendritic}
Bhatt, D.~L., Bhatt, D.~L., and Bhatt, D.~L.\ (2009).
\newblock Dendritic dynamics \emph{in vivo} change during the critical period.
\newblock \emph{Journal of Neurobiology}, 64(1):27--34.

\bibitem[Clark et al.(2019)]{clark2019does}
Clark, K., Khandelwal, U., Levy, O., and Manning, C.~D.\ (2019).
\newblock What does BERT look at? An analysis of BERT's attention.
\newblock In \emph{Proceedings of the 2019 ACL Workshop BlackboxNLP}.

\bibitem[Cobbe et al.(2021)]{cobbe2021gsm8k}
Cobbe, K., Kosaraju, V., Bavarian, M., et al.\ (2021).
\newblock Training verifiers to solve math word problems.
\newblock \emph{arXiv preprint arXiv:2110.14168}.

\bibitem[Dettmers et al.(2023)]{dettmers2023qlora}
Dettmers, T., Pagnoni, A., Holtzman, A., and Zettlemoyer, L.\ (2023).
\newblock QLoRA: Efficient finetuning of quantized LLMs.
\newblock In \emph{Advances in Neural Information Processing Systems (NeurIPS)}.

\bibitem[Fan et al.(2020)]{fan2020reducing}
Fan, A., Grave, E., and Joulin, A.\ (2020).
\newblock Reducing transformer depth on demand with structured dropout.
\newblock In \emph{International Conference on Learning Representations (ICLR)}.

\bibitem[Fedus et al.(2022)]{fedus2022switch}
Fedus, W., Zoph, B., and Shazeer, N.\ (2022).
\newblock Switch transformers: Scaling to trillion parameter models with simple and efficient sparsity.
\newblock \emph{Journal of Machine Learning Research}, 23(120):1--39.

\bibitem[Frankle and Carlin(2019)]{frankle2019lottery}
Frankle, J. and Carlin, M.\ (2019).
\newblock The lottery ticket hypothesis: Finding sparse, trainable neural networks.
\newblock In \emph{International Conference on Learning Representations (ICLR)}.

\bibitem[Gale et al.(2019)]{gale2019state}
Gale, T., Elsen, E., and Hooker, S.\ (2019).
\newblock The state of sparsity in deep neural networks.
\newblock \emph{arXiv preprint arXiv:1902.09574}.

\bibitem[Goyal et al.(2020)]{goyal2020power}
Goyal, S., Choudhury, A., Raje, S., Chakaravarthy, V., Sabharwal, Y., and Verma, A.\ (2020).
\newblock PoWER-BERT: Accelerating BERT inference for classification tasks.
\newblock In \emph{International Conference on Machine Learning (ICML)}.

\bibitem[Graves(2016)]{graves2016adaptive}
Graves, A.\ (2016).
\newblock Adaptive computation time for recurrent neural networks.
\newblock \emph{arXiv preprint arXiv:1603.08983}.

\bibitem[Han et al.(2015)]{han2015learning}
Han, S., Pool, J., Tran, J., and Dally, W.~J.\ (2015).
\newblock Learning both weights and connections for efficient neural networks.
\newblock In \emph{Advances in Neural Information Processing Systems (NeurIPS)}.

\bibitem[Hendrycks et al.(2021)]{hendrycks2021math}
Hendrycks, D., Burns, C., Kadavath, S., et al.\ (2021).
\newblock Measuring mathematical problem solving with the MATH dataset.
\newblock In \emph{Advances in Neural Information Processing Systems (NeurIPS)}.

\bibitem[Hu et al.(2022)]{hu2022lora}
Hu, E.~J., Shen, Y., Wallis, P., et al.\ (2022).
\newblock LoRA: Low-rank adaptation of large language models.
\newblock In \emph{International Conference on Learning Representations (ICLR)}.

\bibitem[Jiang et al.(2023)]{jiang2023mistral}
Jiang, A.~Q., Sablayrolles, A., Mensch, A., et al.\ (2023).
\newblock Mistral 7B.
\newblock \emph{arXiv preprint arXiv:2310.06825}.

\bibitem[Kim and Cho(2020)]{kim2020fastformers}
Kim, Y. and Cho, H.\ (2020).
\newblock Fastformers: Highly efficient transformer models for natural language understanding.
\newblock \emph{arXiv preprint arXiv:2010.13382}.

\bibitem[Li et al.(2017)]{li2017pruning}
Li, H., Kadav, A., Durdanovic, I., Samet, H., and Graf, H.~P.\ (2017).
\newblock Pruning filters for efficient ConvNets.
\newblock In \emph{International Conference on Learning Representations (ICLR)}.

\bibitem[Lightman et al.(2023)]{lightman2023lets}
Lightman, H., Kosaraju, V., Burda, Y., et al.\ (2023).
\newblock Let's verify step by step.
\newblock \emph{arXiv preprint arXiv:2305.20050}.

\bibitem[Loshchilov and Hutter(2019)]{loshchilov2019decoupled}
Loshchilov, I. and Hutter, F.\ (2019).
\newblock Decoupled weight decay regularization.
\newblock In \emph{International Conference on Learning Representations (ICLR)}.

\bibitem[Michel et al.(2019)]{michel2019sixteen}
Michel, P., Levy, O., and Neubig, G.\ (2019).
\newblock Are sixteen heads really better than one?
\newblock In \emph{Advances in Neural Information Processing Systems (NeurIPS)}.

\bibitem[Molchanov et al.(2017)]{molchanov2017variational}
Molchanov, D., Ashukha, A., and Vetrov, D.\ (2017).
\newblock Variational dropout sparsifies deep neural networks.
\newblock In \emph{International Conference on Machine Learning (ICML)}.

\bibitem[Paolicelli et al.(2011)]{paolicelli2011synaptic}
Paolicelli, R.~C. et al.\ (2011).
\newblock Synaptic pruning by microglia is necessary for normal brain development.
\newblock \emph{Science}, 333(6048):1456--1458.

\bibitem[Qwen(2025)]{qwen2025qwen25}
Qwen Team.\ (2025).
\newblock Qwen2.5 technical report.
\newblock \emph{arXiv preprint arXiv:2412.15115}.

\bibitem[Schuster et al.(2022)]{schuster2022confident}
Schuster, T., Fisch, A., Jaakkola, T., and Barzilay, R.\ (2022).
\newblock Confident adaptive language modeling.
\newblock In \emph{Advances in Neural Information Processing Systems (NeurIPS)}.

\bibitem[Shazeer et al.(2017)]{shazeer2017outrageously}
Shazeer, N. et al.\ (2017).
\newblock Outrageously large neural networks: The sparsely-gated mixture-of-experts layer.
\newblock In \emph{International Conference on Learning Representations (ICLR)}.

\bibitem[Stevens et al.(2007)]{stevens2007complement}
Stevens, B. et al.\ (2007).
\newblock The classical complement cascade mediates CNS synapse elimination.
\newblock \emph{Cell}, 131(6):1164--1178.

\bibitem[Suzgun et al.(2022)]{suzgun2022challenging}
Suzgun, M. et al.\ (2022).
\newblock Challenging BIG-bench tasks and whether chain-of-thought can solve them.
\newblock \emph{arXiv preprint arXiv:2210.09261}.

\bibitem[Touvron et al.(2023)]{touvron2023llama}
Touvron, H. et al.\ (2023).
\newblock Llama 2: Open foundation and fine-tuned chat models.
\newblock \emph{arXiv preprint arXiv:2307.09288}.

\bibitem[Vaswani et al.(2017)]{vaswani2017attention}
Vaswani, A., Shazeer, N., Parmar, N., et al.\ (2017).
\newblock Attention is all you need.
\newblock In \emph{Advances in Neural Information Processing Systems (NeurIPS)}.

\bibitem[Voita et al.(2019)]{voita2019analyzing}
Voita, E., Talbot, D., Moiseev, F., Sennrich, R., and Titov, I.\ (2019).
\newblock Analyzing multi-head self-attention: Specialized heads do the heavy lifting, the rest can be pruned.
\newblock In \emph{Proceedings of ACL}.

\bibitem[Wang et al.(2021)]{wang2021spatten}
Wang, H., Zhang, Z., and Han, S.\ (2021).
\newblock SpAtten: Efficient sparse attention architecture with cascade token and head pruning.
\newblock In \emph{IEEE International Symposium on High-Performance Computer Architecture (HPCA)}.

\bibitem[Wei et al.(2022)]{wei2022chain}
Wei, J. et al.\ (2022).
\newblock Chain-of-thought prompting elicits reasoning in large language models.
\newblock In \emph{Advances in Neural Information Processing Systems (NeurIPS)}.

\end{thebibliography}

%==============================================================
% Appendices
%==============================================================
\appendix

\section{Gradient Flow Through Soft Masks}
\label{app:gradients}

We show that the gradient of $\Lcal_{\text{total}}$ with respect to the agent parameters
$\theta_l$ is well-defined for all $T > 0$. Let $z_{l,h} = [f_{\theta_l}(\Phi_l)]_h$ denote
the pre-sigmoid logit. Then $m_{l,h} = \sigma(z_{l,h}/T)$ and:
\[
  \frac{\partial m_{l,h}}{\partial z_{l,h}} = \frac{1}{T}\,m_{l,h}(1 - m_{l,h}),
\]
which is strictly positive for all $m_{l,h} \in (0,1)$. The gradient of the full loss
(Eq.~\ref{eq:loss}) therefore flows cleanly through the masking operation to the agent
parameters throughout training, confirming that the soft-mask formulation is a valid
continuous relaxation of the discrete head-pruning problem.

\section{Entropy Regularization Analysis}
\label{app:entropy}

The entropy term $\Lcal_{\text{entropy}} = H(m_{l,h}) = -m_{l,h}\log m_{l,h}
- (1-m_{l,h})\log(1-m_{l,h})$ is maximized at $m_{l,h} = 0.5$ and equals zero at
$m_{l,h} \in \{0, 1\}$. By \emph{subtracting} this term in Eq.~\eqref{eq:loss},
the optimizer is penalized for maintaining masks near $0.5$ (indecision), and incentivized
to push masks toward the binary extremes. Formally:
\[
  \frac{\partial}{\partial m_{l,h}}(-\beta\Lcal_{\text{entropy}}) =
  \beta\,\log\!\frac{m_{l,h}}{1 - m_{l,h}},
\]
which is negative for $m_{l,h} < 0.5$ (pushing toward 0) and positive for $m_{l,h} > 0.5$
(pushing toward 1), creating a bifurcating force that accelerates convergence to decisive
pruning decisions.

\section{NeurIPS Reproducibility Checklist}
\label{app:reproducibility}

\begin{itemize}
  \item[$\checkmark$] Multiple fixed random seeds used for all experiments
        (Section~\ref{sec:exp_training})
  \item[$\checkmark$] Model and dataset versions specified
        (Sections~\ref{sec:exp_models}--\ref{sec:exp_data})
  \item[$\checkmark$] Fully pinned dependencies: \texttt{pyproject.toml} in repository
  \item[$\checkmark$] Statistical tests described: bootstrap CI, paired bootstrap
        (Section~\ref{sec:exp_eval})
  \item[$\checkmark$] Code available:
        \url{https://github.com/Tommaso-R-Marena/microglia-pruning}
  \item[$\checkmark$] All datasets publicly available: GSM8K, MATH, BIG-Bench
\end{itemize}

\end{document}